% TDP-Reasoner paper draft. NeurIPS-like single column (5.5 in text width, Times). Sections not yet written are marked.
\documentclass[10pt]{article}
\usepackage[letterpaper,textwidth=5.5in,textheight=9in,top=1in]{geometry}
\usepackage[T1]{fontenc}
\usepackage{mathptmx}
\usepackage{amsmath,amssymb}
\usepackage{microtype}
\usepackage[hidelinks]{hyperref}

\setlength{\parindent}{0pt}
\setlength{\parskip}{5pt}
\newcommand{\todo}[1]{\textit{[#1]}}
\newcommand{\R}{\mathbb{R}}
\newcommand{\kt}{\tilde k}
\newcommand{\qt}{\tilde q}
\renewcommand{\Re}{\operatorname{Re}}
\DeclareMathOperator{\RMSNorm}{RMSNorm}
\DeclareMathOperator{\RoPE}{RoPE}
\DeclareMathOperator{\FFN}{FFN}
\DeclareMathOperator{\soft}{soft}

\begin{document}

\vbox{%
  \hrule height 4pt
  \vskip 0.25in
  \centering
  {\LARGE\bfseries TDP-Reasoner: Timing-Dependent Plasticity\\[2pt] Along the Recursion\par}
  \vskip 0.29in
  \hrule height 1pt
}
\vskip 0.3in
\begin{center}
  {\bfseries Anonymous Author(s)}\\
  Affiliation\\
  \texttt{email}
\end{center}
\vskip 0.3in

\begin{center}
  {\large\bfseries Abstract}
\end{center}
\vskip 0.05in
\begin{list}{}{\leftmargin=0.5in \rightmargin=0.5in}\item[]
Small recursive models trained from scratch are the strongest non-LLM approach to puzzle benchmarks such as
Sudoku-Extreme and ARC-AGI. These models, HRM and TRM among them, re-apply a Transformer block many times. Seen as a
modern Hopfield network, each step's attention is a Hebbian memory of the current state, so recursion reinforces the
model's current answer, and accuracy falls once it runs past its training horizon. We propose TDP-Reasoner
(timing-dependent plasticity reasoner), which keeps a fast-weight memory along the recursion itself: keys and values from
different steps are paired under a learned two-sided timing window, borrowed from spike-timing-dependent plasticity, so
the memory tracks how the state moves rather than where it is. Demonstrations of a task reach the query only through a
per-task linear rule fitted in context, so unseen tasks need no retraining. Trained for 256 recursion blocks and run for
10{,}000, it raises held-out Sudoku-Extreme accuracy from 59.8\% to 82.0\% and never loses a solved puzzle. On unseen
ARC-AGI-1 tasks with small grids it matches the public TRM checkpoint (43.2\% vs.\ 45.5\%) after about 3\% of its
training; at full scale it reaches 18.8\% pass@2 on ARC-AGI-1 (TRM: 41.8\% under the same evaluation) and 0\% on
ARC-AGI-2.
\end{list}

\section{Introduction}
\todo{To be written.}

\section{Related work}
\todo{To be written.}

\section{TDP-Reasoner}
\label{sec:method}

We start from recurrent linear attention, add a memory across the recursion (the Sudoku model), and then add a rule
head that lets the model read the demonstrations of an ARC task.

\subsection{Starting point: recurrent linear attention}
\label{sec:block}

One block with shared weights is applied at steps $n = 1, 2, \dots$ of the recursion. At step $n$ it reads
\begin{equation}
  \operatorname{read}_n = \Big(\textstyle\sum_{\tau} v_{n,\tau}\, k_{n,\tau}^{\top}\Big)\, q_n ,
  \label{eq:linattn}
\end{equation}
linear attention over the tokens $\tau$ with unit-length queries and keys (QK-norm), followed by the usual residual,
RMSNorm and feed-forward layer (Appendix~\ref{app:block}). Within a step this is ordinary attention. Across steps
nothing is kept except the state itself.

\subsection{Traces and a timing window}
\label{sec:memory}

TDP-Reasoner adds a fast weight $W$ that is carried from step to step:
\begin{equation}
  \operatorname{read}_n = \Big(\textstyle\sum_{\tau} v_{n,\tau}\, k_{n,\tau}^{\top} + W_n\Big)\, q_n .
  \label{eq:read}
\end{equation}
$W_n$ holds every pair of a value and a key of the same token taken at two steps $t, s \le n$, weighted by a window of
their time difference:
\begin{equation}
  W_n = \sum_{\tau} \sum_{s,\,t \,\le\, n} \rho_{\,n-\max(s,t)}\; L(t-s)\; v_{t,\tau}\, k_{s,\tau}^{\top}.
  \label{eq:pairs}
\end{equation}
With the key as presynaptic and the value as postsynaptic activity, this is the pair-based STDP integral
$\Delta w = \sum W(t_{\mathrm{post}} - t_{\mathrm{pre}})\, \mathrm{post}\cdot\mathrm{pre}$: pairs in which the key came
first ($t > s$) fall on one side of the window and pairs in which the value came first on the other, so $L$ can
strengthen one order and weaken the other. The factor $\rho_a = \eta(1-\eta)^a$ slowly forgets old pairs and keeps $W$
bounded. Equation~\eqref{eq:pairs} is computed online with two decaying traces per token, one of keys and one of
values, at the cost of two matrix products per step; $L$ is a sum of a few damped complex exponentials whose poles and
coefficients are learned (Appendix~\ref{app:memory}; Section~\ref{sec:theory} proves that the online computation is
exact). \todo{Figure: windows learned on Sudoku.}

This is the Sudoku model: recurrent linear attention with QK-norm, plus traces and a timing window.

\subsection{Reading demonstrations}
\label{sec:rule}

For ARC every example, an input grid with its output grid, is a \emph{lane} that the block processes on its own; the
output grid is hidden in every lane, demonstration or query. The lanes of one task form a \emph{group}. After every
block a rule head fits, for each lane, a linear map $\Theta$ from cell features $\phi$ to colours over the \emph{other}
demonstrations of its group, one lasso step at a time:
\begin{equation}
  \Theta \leftarrow \operatorname{soft}_\lambda\!\big(\Theta - (\Theta C - B)\big), \qquad r = \Theta\,\phi ,
  \label{eq:rule}
\end{equation}
where $C$ and $B$ average $\phi\phi^{\top}$ and $y\,\phi^{\top}$ ($y$ the one-hot colour) over those demonstrations'
output cells, and $\operatorname{soft}_\lambda$ is soft-thresholding. The prediction $r$ is added to the next block's
input on the output canvas. A query never enters $C$ or $B$, so it is answered from the demonstrations alone, with no
puzzle embedding and no test-time training. As in the Tolman--Eichenbaum machine, the features are shared structure and
$\Theta$ binds them to content for one task; a few position frames in the rule head retrieve colours by location
(Appendix~\ref{app:rule}).

\subsection{Training and inference}
\label{sec:training}

A segment runs a fixed number of blocks and backpropagates through the last few; the state, traces, $W$ and $\Theta$
are carried to the next segment without gradient, with a loss after every segment (TRM's deep supervision without
halting) and TRM's stablemax cross-entropy. At inference, as in TRM, each task is run in several augmented versions and
the answers are voted; ties are broken by how well each version predicts its own demonstrations
(Appendix~\ref{app:block}).

\section{Theory}
\label{sec:theory}
\todo{To be written: the writes equal the pair-based STDP integral; the window sees only the lagged cross-correlation of
activities; recursion supplies temporal order.}

\section{Experiments}
\label{sec:experiments}
\todo{To be written.}

\section{Limitations and future work}
\todo{To be written.}

\appendix

\section{The block, segments and inference}
\label{app:block}

A lane's state $h \in \R^{T\times d}$ starts from a fixed vector. At step $n$, with $x$ the lane's tokens and
$r_{n-1}$ the rule head's prediction on the output canvas,
\begin{align*}
  u_n &= h_{n-1} + \sqrt{d}\, e_n, \qquad
  e_n = \operatorname{Embed}(x) + \begin{bmatrix} 0 \\ W_{\mathrm{fb}}\, r_{n-1} \end{bmatrix}, \\
  \qt_n &= \RoPE\!\big(q_n/\lVert q_n\rVert\big), \quad
  \kt_n = \RoPE\!\big(k_n/\lVert k_n\rVert\big), \quad
  [q_n; k_n; v_n] = [W_q; W_k; W_v]\, u_n \quad \text{(per head)}, \\
  h'_n &= \RMSNorm\big(u_n + W_O \operatorname{read}_n\big), \qquad
  h_n = \RMSNorm\big(h'_n + \FFN(h'_n)\big),
\end{align*}
with $\operatorname{read}_n$ from \eqref{eq:read} computed per head, $\FFN$ a SwiGLU layer and parameter-free
$\RMSNorm$. $\RoPE$ uses either fixed one-dimensional frequencies over the token index or learned angles on canvas rows
and columns (Section~\ref{sec:experiments} says which run uses which). An ARC lane stacks the input and output canvases
($C \times C$ each, $C = 30$) into $T = 2C^2$ tokens over \{unknown, pad, colours 0--9, end of grid\}.

\todo{Segments: blocks per segment and backpropagated blocks per run; 16 segments per lane in full-scale runs, one in
the small-grid harness. Inference: versions, inverse transforms, voting, tie-break by demonstration consistency.}

\section{Computing the memory online}
\label{app:memory}

For every head and token, each band $r = 1,\dots,R$ keeps a key trace and a value trace over the steps:
\begin{align*}
  E^{K}_{r}(n) &= z_r\, E^{K}_{r}(n-1) + \gamma_r\, \kt_n, \qquad
  D^{V}_{r}(n) = z_r\, E^{V}_{r}(n-1), \qquad E^{V}_{r}(n) = D^{V}_{r}(n) + \gamma_r\, v_n, \\
  P_n &= \sum_{\tau} \sum_{r=1}^{R} \Re\!\Big( a_r\, v_n\, E^{K}_{r}(n)^{\top} + b_r\, D^{V}_{r}(n)\, \kt_n^{\top} \Big),
  \qquad W_n = (1-\eta)\, W_{n-1} + \eta\, P_n ,
\end{align*}
with learned poles $z_r = \exp(-e^{\nu_r} + i\theta_r)$, $\gamma_r = \sqrt{1 - |z_r|^2}$ (the parameterisation of
linear recurrent units with $\gamma_r$ tied to the pole), learned complex coefficients $a_r, b_r$ and a learned rate
$\eta \le 1/8$, all per head. The window of \eqref{eq:pairs} is
\[
  L(\Delta) = \sum_r \gamma_r \Re\big(a_r z_r^{\Delta}\big) \ (\Delta \ge 0), \qquad
  L(\Delta) = \sum_r \gamma_r \Re\big(b_r z_r^{|\Delta|}\big) \ (\Delta < 0),
\]
which contains the usual exponential STDP window ($R = 1$, $\theta = 0$, $a > 0$, $b < 0$). We use $R = 4$ bands with
initial periods of 2, 6, 20 and 64 steps; $a_r, b_r$ start small and random. The bands are summed before the token sums,
so a step costs two matrix products for any $R$. Traces and $W$ start at zero for a new lane.

\section{Rule head}
\label{app:rule}

Features are $\phi(o) = W_\phi h(o) / \lVert W_\phi h(o) \rVert \in \R^{p}$ on the output-canvas cells $o$. For a lane
$\ell$, with $n_\ell$ the number of labelled output cells in the other demonstration lanes $\ell'$ of its group,
\[
  C_\ell = \frac{1}{n_\ell} \sum_{\ell' \ne \ell} \sum_{o} \phi_{\ell'}(o)\, \phi_{\ell'}(o)^{\top}, \qquad
  B_\ell = \frac{1}{n_\ell} \sum_{\ell' \ne \ell} \sum_{o} y_{\ell'}(o)\, \phi_{\ell'}(o)^{\top},
\]
and \eqref{eq:rule} is one proximal step on
$\tfrac12 \operatorname{tr}(\Theta C_\ell \Theta^{\top}) - \operatorname{tr}(\Theta B_\ell^{\top}) + \lambda \lVert \Theta \rVert_1$,
valid at unit step because $\lambda_{\max}(C_\ell) \le \operatorname{tr} C_\ell = 1$. The prediction has a frame term,
$r(o) = \Theta \phi(o) + \sum_{k=0}^{K} \beta_k \kappa_k(o)$: $\kappa_0$ is the colour histogram of the lane's input
grid and $\kappa_k$ ($k \ge 1$) a softmax read of the input grid's colour cells from cell $o$ with learned rotary
frequencies on canvas coordinates ($K = 8$, scaled by $1/\sqrt{K+1}$); $\beta$ belongs to the same lasso. $r$ enters the
next block through $W_{\mathrm{fb}}$, initialised at zero. The logits are
$\operatorname{classifier}(h(o)) + \tau\, r(o)$, or an output FFN applied to $h(o) + \sqrt{d}\, W_{\mathrm{fb}} r(o)$
(Section~\ref{sec:experiments}).

\end{document}
